\chapter{Appendice}
\section{Metodo dei minimi quadrati per il calcolo stabile della derivata}
Di seguito viene illustrato il procedimento matematico effettuato per ottenere la formula utilizzata nel controllore PID per
il calcolo del contributo derivativo.

Il calcolo della derivata di una grandezza fisica misurata da un sensore risulta particolarmente delicato in quanto è fortemente
soggetto al rumore delle misurazioni. Il semplice calcolo del rapporto incrementale tra due campioni consecutivi può portare
a rapide variazioni della derivata che potrebbero indurre problemi di stabilità del controllore. Per mitigare le fluttuazioni
istantanee si è scelto di utilizzare il metodo dei minimi quadrati ed effettuare una regressione lineare su un set più ampio
di dati in modo da compensare le fluttuazioni istantanee e ottenere una stima più stabile della derivata.

\subsection*{Notazioni e convenzioni utilizzate}
Nei calcoli successivi si adottano le seguenti notazioni:
\begin{itemize}
	\item \(y_i\) : misurazione del sensore all'istante di tempo \(x_i\)
	\item \(a_0 + a_1 \, x_i\) : retta di regressione lineare, dove \(a_0\) è intercetta e \(a_1\) è il coefficiente angolare
	della retta di regressione lineare, che rappresenta la derivata da calcolare
\end{itemize}

Per calcolare la derivata ad un certo istante di tempo \(x_i\), si considerano solo gli ultimi \(N\) campioni misurati dal
sensore. Scelto \(x_k\) l'istante di tempo associato all'ultima misurazione \(y_j\), l'indice \(i\) dei campioni considerati
varia in \([j \!-\! (N \!-\! 1), \, j]\) con \(j \in \mathbb{N}\). Siccome gli indici delle sommatorie sono variabili mute,
nei successivi passaggi si considera senza perdita di generalità \(j = N \!-\! 1\) in modo che \(i\) varia da \(0\) a \(N-1\).

\subsection*{Procedimento matematico}
L'errore \(e_i\) tra la misurazione \(y_i\) e la retta di regressione lineare \(a_0 + a_1 \, x_i\) viene definito come:
\begin{equation}
	e_i = y_i - (a_0 + a_1 \, x_i)
\end{equation}

Si definisce quindi l'errore quadratico totale \(S(a_0, a_1)\) come segue:
\begin{equation}
	S(a_0, a_1) \; = \; \sum_{i = 0}^{N-1} e_i^2 \; = \; \sum_{i = 0}^{N-1} (y_i - (a_0 + a_1 \, x_i))^2
\end{equation}

I parametri \(a_0\) e \(a_1\) della retta di regressione lineare che minimizzano l'errore quadratico totale \(S(a_0, a_1)\)
si ottengono ponendo a zero le derivate parziali di \(S(a_0, a_1)\) rispetto a \(a_0\) e \(a_1\) e risolvendo il sistema
lineare di due equazioni a due incognite \(a_0\) e \(a_1\) così ottenuto.
\begin{align}
	\frac{\partial S(a_0, a_1)}{\partial a_0} &= -2 \sum_{i = 0}^{N-1} (y_i - (a_0 + a_1 \, x_i)) = \\
	&= -2 \left( \sum_{i = 0}^{N-1} y_i - \sum_{i=0}^{N-1} a_0 - \sum_{i=0}^{N-1} a_1 \, x_i \right) = \\
	&= -2 \left( \sum_{i = 0}^{N-1} y_i - a_0 \sum_{i=0}^{N-1} 1 - a_1 \sum_{i=0}^{N-1} x_i \right) = \\
	&= -2 \left( N \bar{y} - a_0 N - a_1 N \bar{x} \right) = \\
	&= -2 N \left( \bar{y} - a_0 - a_1 \bar{x} \right)
\end{align}
\begin{align}
	\frac{\partial S(a_0, a_1)}{\partial a_1} &= -2 \sum_{i = 0}^{N-1} x_i (y_i - (a_0 + a_1 \, x_i)) = \\
	&= -2 \left( \sum_{i = 0}^{N-1} x_i y_i - \sum_{i=0}^{N-1} a_0 x_i - \sum_{i=0}^{N-1} a_1 {x_i}^2 \right) = \\
	&= -2 \left( \sum_{i = 0}^{N-1} x_i y_i - a_0 \sum_{i=0}^{N-1} x_i - a_1 \sum_{i=0}^{N-1} {x_i}^2 \right) = \\
	&= -2 \left( \sum_{i = 0}^{N-1} x_i y_i - a_0 N \bar{x} - a_1 \sum_{i=0}^{N-1} {x_i}^2 \right)
\end{align}

Si procede quindi a risolvere il sistema di equazioni lineari così ottenuto per sostituzione:
\begin{equation}
	\begin{cases}
		\frac{\partial S(a_0, a_1)}{\partial a_0} = 0 \\
		\frac{\partial S(a_0, a_1)}{\partial a_1} = 0
	\end{cases}
	\quad \Rightarrow \quad
	\begin{cases}
		2 N \left( \bar{y} - a_0 - a_1 \bar{x} \right) = 0 \\
		-2 \left( \sum_{i = 0}^{N-1} x_i y_i - a_0 N \bar{x} - a_1 \sum_{i=0}^{N-1} {x_i}^2 \right) = 0
	\end{cases}
\end{equation}

\begin{equation}
	-2 N \left( \bar{y} - a_0 - a_1 \bar{x} \right) = 0 \quad \Rightarrow \quad a_0 = \bar{y} - a_1 \bar{x}
\end{equation}
\begin{align}
	-2 \left( \sum_{i = 0}^{N-1} x_i y_i - a_0 N \bar{x} - a_1 \sum_{i=0}^{N-1} {x_i}^2 \right) = 0 \\
	\sum_{i = 0}^{N-1} x_i y_i - (\bar{y} - a_1 \bar{x}) N \bar{x} - a_1 \sum_{i=0}^{N-1} {x_i}^2 = 0 \\
	\sum_{i = 0}^{N-1} x_i y_i - N \bar{y} \bar{x} + a_1 N {\bar{x}}^2 - a_1 \sum_{i=0}^{N-1} {x_i}^2 = 0 \\
	\sum_{i = 0}^{N-1} x_i y_i - N \bar{y} \bar{x} - a_1 \left( \sum_{i=0}^{N-1} {x_i}^2 - N {\bar{x}}^2 \right) = 0 \\
\end{align}
\begin{equation}
	a_1 = \frac{\displaystyle \;\; \sum_{i=0}^{N-1} x_i y_i - N \bar{y} \bar{x} \;\;}{\displaystyle \;\; \sum_{i=0}^{N-1} {x_i}^2 - N {\bar{x}}^2 \;\;} \label{eq:7_a1_base}
\end{equation}

Si assume che gli istanti di tempo \(x_i\) sono tutti equidistanti, ovvero che formano una sequenza continua dei multipli del
tempo di campionamento \(T_s\) a partire dall'istante iniziale \(k \, T_s\) dove \(k = j - (N - 1)\). I campioni \(x_i\) possono
essere quindi espressi in funzione dei due indici \(i\) e \(k\):
\begin{equation}
	x_i = (j - (N-1)) \, T_s + i \, T_s = (k + i) \, T_s \qquad \text{con} \; \bar{x} = \frac{1}{N} \sum_{i=0}^{N-1} x_i = \frac{T_s}{N} \sum_{i=0}^{N-1} (k + i)\label{eq:7_sub}
\end{equation}

È possibile quindi riscrivere la formula \eqref{eq:7_a1_base} con la sostituzione \eqref{eq:7_sub} individuata sopra:
\begin{equation}
	a_1 = \frac{\displaystyle \;\; T_s \sum_{i=0}^{N-1} (k + i) \, y_i - N \bar{y} \left(\frac{T_s}{N} \, \sum_{i=0}^{N-1} (k + i)\right) \;\;}{\displaystyle \;\; {T_s}^2 \sum_{i=0}^{N-1} (k + i)^2 - N \left(\frac{T_s}{N} \,\sum_{i=0}^{N-1} (k + i)\right)^2 \;\;}
\end{equation}

\newpage

Analizzando separatamente il numeratore e il denominatore della formula \eqref{eq:7_a1_base} si ottengono le seguenti
semplificazioni fino ad ottenere la formula finale utilizzata nell'implementazione del controllore PID:
\begin{align}
	\textit{Num.} &= T_s \sum_{i=0}^{N-1} (k + i) \, y_i - N \bar{y} \left(\frac{T_s}{N} \, \sum_{i=0}^{N-1} (k + i)\right) = \\
	&= T_s \left(k \sum_{i=0}^{N-1} y_i + \sum_{i=0}^{N-1} i \cdot y_i\right) - N \bar{y} \cdot \frac{T_s}{N} \cdot \left( \sum_{i=0}^{N-1} i + \sum_{i=0}^{N-1} k \right) = \\
	&= T_s \left( k N \bar{y} + \sum_{i=0}^{N-1} i \cdot y_i - \bar{y} \left(\frac{N(N-1)}{2} + kN\right) \right) = \\
	&= T_s \left( k N \bar{y} + \sum_{i=0}^{N-1} i \cdot y_i - \frac{N(N-1)}{2} \cdot \bar{y} - k N \bar{y} \right) = \\
	&= T_s \left( \sum_{i=0}^{N-1} i \cdot y_i - \frac{N(N-1)}{2} \cdot \bar{y} \right) = \\
	&= T_s \left( \sum_{i=0}^{N-1} i \cdot y_i - \frac{N-1}{2} \cdot \sum_{i=0}^{N-1} y_i \right) \label{eq:7_num}
\end{align}
\begin{align}
	\textit{Den.} &= {T_s}^2 \sum_{i=0}^{N-1} (k + i)^2 - N \left(\frac{T_s}{N} \,\sum_{i=0}^{N-1} (k + i)\right)^2 = \\
	&= {T_s}^2 \left( \sum_{i=0}^{N-1} (k^2 + 2ik + i^2) - \frac{T_s^2}{N} \left( \sum_{i=0}^{N-1} (k + i) \right)^2 \right) = \\
	&= {T_s}^2 \left( k^2 N + 2k \sum_{i=0}^{N-1} i + \sum_{i=0}^{N-1} i^2 - \frac{1}{N} \left( k N + \sum_{i=0}^{N-1} i \right)^2 \right) = \\
	&= {T_s}^2 \left( k^2 N \!+\! 2k \frac{N(N\!\!-\!\!1)}{2} \!+\! \frac{N(N\!\!-\!\!1)(2N\!\!-\!\!1)}{6} - \frac{1}{N} \left( k N \!+\! \frac{N(N\!\!-\!\!1)}{2} \right)^2 \right) = \\
	\begin{split}
		&= {T_s}^2 \bigg( k^2 N + k N^2 - kN + \frac{2N^3 - 3 N^2 + N}{6} + \\[-5pt]
		&\qquad\qquad\qquad\qquad -\frac{1}{N} \cdot \bigg( k^2 N^2 + kN^3 - kN^2 + \frac{N^4 - 2N^3 + N^2}{4} \bigg) \bigg) = \\
	\end{split} \\
	\begin{split}
		&= \frac{{T_s}^2}{12} \cdot \big( 12k^2 N + 12k N^2 - 12kN + 4N^3 - 6N^2 + 2N + \\
		&\qquad\qquad\qquad\qquad + 12k^2 N - 12k N^2 + 12kN - 3N^3 + 6N^2 - 3N \big) =
	\end{split} \\
	&= \frac{{T_s}^2}{12} \left( N^3 - N \right) = \frac{{T_s}^2}{12} \cdot N (N-1)(N+1) \label{eq:7_den}
\end{align}


Unendo i due termini \eqref{eq:7_num} e \eqref{eq:7_den} così ottenuti, si arriva alla formula finale utilizzata
nell'implementazione del controllore PID.
\begin{align}
	a_1 &= \frac{\displaystyle \;\; T_s \sum_{i=0}^{N-1} (k + i) \, y_i - N \bar{y} \left(\frac{T_s}{N} \, \sum_{i=0}^{N-1} (k + i)\right) \;\;}{\displaystyle \;\; {T_s}^2 \sum_{i=0}^{N-1} (k + i)^2 - N \left(\frac{T_s}{N} \,\sum_{i=0}^{N-1} (k + i)\right)^2 \;\;} = \\[5pt]
	&= \frac{\displaystyle \;\; T_s \left( \sum_{i=0}^{N-1} i \cdot y_i - \frac{N-1}{2} \cdot \sum_{i=0}^{N-1} y_i \right) \;\;}{\displaystyle \;\; \frac{{T_s}^2}{12} \cdot N (N-1)(N+1) \;\;} = \\[5pt]
	&= \frac{12}{T_s \cdot N(N-1)(N+1)} \cdot \left( \sum_{i=0}^{N-1} i \cdot y_i - \frac{N-1}{2} \cdot \sum_{i=0}^{N-1} y_i \right) = \\
	&= \frac{12}{T_s \cdot N(N^2-1)} \cdot \sum_{i=0}^{N-1} i \cdot y_i - \frac{6}{T_s \cdot N(N+1)} \cdot \sum_{i=0}^{N-1} y_i
\end{align}
